\documentclass[preprint,authoryear,review,12pt]{elsarticle}
\usepackage{amsmath}
\usepackage{amsfonts}
\usepackage{lscape}
\usepackage{fullpage}
\usepackage{graphicx,color} % Para escrever em cores
\usepackage{graphics}
\usepackage{subfig}
\usepackage{ae}
\usepackage[T1]{fontenc}
\usepackage[utf8]{inputenc}	
%\usepackage[ansinew]{inputenc}
\usepackage{pdflscape}
\usepackage{url}
\usepackage{setspace}
\usepackage{tabularx}
\usepackage{dcolumn}
\usepackage{rotating}
\usepackage{a4wide}
\usepackage{supertabular}
\usepackage[lmargin=1.5cm, rmargin=1.5cm,tmargin=1.50cm,bmargin=1.50cm]{geometry}
\usepackage{fullpage}
\usepackage{multirow}
\usepackage{longtable}
\usepackage{geometry}
\usepackage[all]{xy}
\usepackage{lscape}
\usepackage{pdflscape}
\usepackage{tabularx}
\usepackage{tabulary}
\usepackage{textgreek}
\usepackage{fixmath}
\usepackage{longtable}
\usepackage{subfig}
\usepackage{graphicx}

\usepackage{booktabs}         % Recommended for professional cell borders
\usepackage[flushleft]{threeparttable}% [flushleft] left-aligns the footnotes
\usepackage{threeparttablex}
\usepackage{longtable}

\usepackage{nicematrix}
\usepackage{enumitem}

%
\usepackage[table,xcdraw]{xcolor}
\usepackage{lipsum}
%\usepackage[margin=1in]{geometry} 
\usepackage{mathtools}
\usepackage{upgreek}
\usepackage{booktabs,caption} %
\usepackage[flushleft]{threeparttable} % Rodapé
\usepackage{float} % Fixar local tabela [H]
\usepackage{arydshln} 
\usepackage[normalem]{ulem}
 \useunder{\uline}{\ul}{}
%
\usepackage{calrsfs}
\definecolor{mygray}{gray}{0.6}
\definecolor{myblue}{RGB}{30 0, 180}

% The amssymb package provides various useful mathematical symbols
\usepackage{mathtools, amsmath, amssymb, amsthm, latexsym}

%\usepackage[top=2cm, left=2cm, right=2cm, bottom=2cm]{geometry}
\usepackage{rotating}

\newcommand{\m}[3]{\multicolumn{#1}{#2}{#3}}
\newcolumntype{d}[1]{D{.}{.}{#1}}

\def\mathbi#1{\textbf{\em #1}}
\newcommand{\mat}[1]{\mbox{\boldmath{$#1$}}}

%%% Nota: abrir com Western (ISO Latin 1) %%%%%
\usepackage{hyperref} % para as citaes e refs voltarem a cor normal comentar esta parte
\hypersetup{
  colorlinks = true,
  linkcolor=blue,   % color of internal links
  citecolor=blue,   % color of links to bibliography
  urlcolor=blue,    % color of external links
  pagebackref=true,
  implicit=false,
  bookmarks=true,
  bookmarksopen=true,
  pdfdisplaydoctitle=true
} 

%\usepackage{dcolumn}

%\journal{Empirical Economics}
\makeatletter
\def\ps@pprintTitle{%
 \let\@oddhead\@empty
 \let\@evenhead\@empty
 \def\@oddfoot{}%
 \let\@evenfoot\@oddfoot}
\makeatother

\begin{document}

\begin{frontmatter}
\title{Uncovered Interest Parity, Inverted: Risk-Premium Transmission of Monetary Policy Shocks in Brazil}

% Uncovered Interest Parity, Inverted: Risk-Premium Transmission of Monetary Policy Shocks in Brazil
% Beyond Uncovered Interest Parity: Risk-Premium Transmission of Monetary Policy Shocks in Brazil


\author[Label0]{Gabriel Arruda}
\author[Label1]{Jo\~ao F. Caldeira }
%\author[Label1]{Guilherme V. Moura}
\address[Label0]{Graduate School of Economics\\Universidade Federal de Santa Catarina}
\address[Label1]{Department of Economics\\Universidade de Brasília \& CNPq}

\begin{abstract}
Ao contrário do que prevê o mecanismo tradicional da paridade descoberta da taxa de juros, um choque contracionista de 50 pontos-base na política monetária brasileira deprecia a moeda doméstica e eleva o prêmio de risco soberano. Após o choque, a taxa de câmbio BRL/USD deprecia $5,\!55\%$ no impacto, enquanto o EMBI+ abre 32,0 pontos-base e o CDS soberano de cinco anos abre 43,4 pontos-base. As três bandas de 90\% excluem zero, de modo que a resposta conjunta é compatível com um canal de transmissão via prêmio de risco. A curva de juros acompanha esse movimento, pois a resposta cresce com a maturidade até cinco anos e alcança 117,0 pontos-base nesse vértice. Esses resultados contrastam com estimativas diárias por identificação via heterocedasticidade que rejeitam a hipótese de dominância fiscal para o Brasil. Estimamos um Modelo de Fatores Dinâmicos Estruturais identificado por instrumento externo, combinando surpresas monetárias em reuniões do Copom, ortogonalização em relação à informação pública predeterminada e um filtro de sinal, sobre 106 séries mensais entre 2013 e 2025.\\

\noindent Abstract: Contrary to the conventional prediction of Uncovered Interest Parity, a 50-basis-point contractionary monetary policy shock in Brazil depreciates the BRL/USD exchange rate by $5.55\%$ on impact. The EMBI+ and the five-year sovereign CDS spread rise by 32.0 and 43.4 basis points, respectively, and all three 90\% confidence bands exclude zero. This joint response is consistent with a risk-premium transmission channel. The yield curve moves in the same direction, with the response increasing by maturity up to 117.0 basis points at five years. These estimates differ from recent daily evidence based on heteroskedasticity identification that rejects fiscal dominance in Brazil. We estimate a Structural Dynamic Factor Model identified through an external instrument that combines Copom surprises, orthogonalization with respect to predetermined public information, and a sign filter, using 106 monthly Brazilian macroeconomic and financial series from 2013 to 2025.

\end{abstract}



\begin{keyword}
Structural Factor Models  \sep Monetary Policy  \sep Asset Prices \sep Sovereign Risk \sep External Instruments
\\[0.1in]
JEL C32 \sep E52 \sep E43 \\[0.1in]
\end{keyword}


\end{frontmatter}
%\newpage
 %\singlespacing  
\onehalfspacing   
% \doublespacing  
%
% henceforth DNS
\section{Introduction}
A compreensão dos canais de transmissão da política monetária constitui um dos temas centrais da macroeconomia, uma vez que as decisões do banco central afetam diretamente as condições de crédito, as expectativas dos agentes e a atividade econômica. Entre esses canais, o canal cambial ocupa posição de destaque. Segundo o mecanismo convencional da paridade descoberta da taxa de juros, um aperto monetário amplia o diferencial de juros doméstico, atrai capital estrangeiro e aprecia a moeda doméstica. Essa previsão, entretanto, pressupõe que o diferencial de juros seja a principal margem de ajuste dos fluxos internacionais de capitais, hipótese que pode deixar de ser válida quando preocupações com a sustentabilidade fiscal passam a desempenhar papel relevante na precificação dos ativos. Esse tipo de reversão da previsão convencional da paridade descoberta não é exclusivo de contextos de dominância fiscal: \cite{dalgicozhan2026} mostram, em um arcabouço teórico e empírico para economias emergentes, que a interação entre descasamento cambial de passivos bancários e faturamento de exportações em dólar (dominant currency pricing) também é capaz de gerar prêmios de risco cambial endógenos que dominam o canal convencional do diferencial de juros, fazendo com que depreciações ocorram justamente nos estados em que os ativos domésticos deveriam, pela UIP padrão, se valorizar.

%A compreensão dos canais de transmissão da política monetária é um dos temas centrais da macroeconomia, pois as decisões do banco central afetam diretamente as condições de crédito, as expectativas dos agentes e a atividade econômica. Entre os canais de transmissão, a resposta da taxa de câmbio ocupa um lugar particular: segundo o mecanismo convencional da paridade descoberta da taxa de juros, um aperto monetário deveria atrair capital estrangeiro e apreciar a moeda doméstica. Esse mecanismo, no entanto, pressupõe que o diferencial de juros seja a única margem relevante de ajuste, uma suposição que pode falhar em economias cuja trajetória fiscal é vista pelos mercados como uma restrição ativa.

Um aperto monetário de 50 pontos-base deprecia o real em $5{,}55\%$, abre o EMBI+ em 32,0 pontos-base e eleva o CDS soberano de cinco anos em 43,4 pontos-base no impacto. As três bandas de 90\% excluem zero, resposta conjunta difícil de reconciliar com a paridade descoberta da taxa de juros e compatível com transmissão via prêmio de risco. A curva acompanha esse movimento, pois o repasse aumenta com a maturidade até 117,0 pontos-base no vértice de cinco anos. Esses resultados são compatíveis com uma leitura de dominância fiscal, embora o desenho não identifique diretamente a origem fiscal do prêmio de risco e a força do instrumento na amostra completa, $\xi_{mp}=7{,}65$, exija cautela na inferência.

%Neste artigo, documentamos um fato estilizado que contraria essa previsão: em determinadas economias emergentes, um choque contracionista de política monetária pode elevar o prêmio de risco soberano a ponto de inverter a resposta esperada da taxa de câmbio. Usando dados brasileiros entre 2013 e 2025, mostramos que um aperto monetário de 50 pontos-base deprecia o real e, simultaneamente, eleva o EMBI+ e o CDS soberano de cinco anos, uma resposta conjunta incompatível com o canal tradicional da paridade de juros, mas consistente com um regime em que o aumento da taxa de juros piora a percepção sobre a sustentabilidade da dívida e esse efeito domina o diferencial de juros convencional. A curva de juros reforça essa leitura: o repasse do choque cresce com a maturidade até cinco anos, superando a amplificação documentada em desenhos comparáveis para economias desenvolvidas, e reverte no médio prazo em sincronia com o câmbio e o prêmio de risco soberano.

Até onde sabemos, este é o primeiro trabalho a documentar esse mecanismo de transmissão para a economia brasileira utilizando um Modelo de Fatores Dinâmicos Estruturais identificado por instrumento externo. Ao relacionar conjuntamente a resposta da taxa de câmbio, do prêmio de risco soberano e da curva de juros, o artigo fornece evidência de que o canal de prêmio de risco pode desempenhar papel relevante na transmissão da política monetária em economias emergentes com fragilidades fiscais como a brasileira.

Identificar empiricamente esse mecanismo de transmissão, entretanto, não é uma tarefa trivial. Choques de política monetária não são diretamente observáveis e sua identificação depende da informação disponível ao econometrista. Modelos tradicionais, como os Vetores Autorregressivos Estruturais (SVAR), frequentemente sofrem do problema da não-fundamentalidade, uma vez que operam com conjuntos de informação muito menores do que aqueles observados pelos agentes econômicos e pelo próprio banco central. Como consequência, podem subestimar a magnitude dos choques monetários ou produzir respostas aparentemente contraintuitivas para algumas variáveis macroeconômicas e financeiras \citep[]{villaverde, alessi}. Para contornar essa limitação, empregamos um Modelo de Fatores Dinâmicos Estruturais (SDFM), que extrai fatores comuns de um amplo conjunto de variáveis macroeconômicas e financeiras, aproximando o conjunto de informação relevante utilizado pelos agentes econômicos e permitindo uma identificação mais robusta dos choques monetários \citep[]{STOCK2016415, alessi}.

%A identificação correta desse mecanismo depende crucialmente da informação disponível para o econometrista. Modelos tradicionais, como os Vetores Autorregressivos Estruturais (SVAR), frequentemente sofrem do problema da não-fundamentalidade, pois não conseguem incorporar o vasto conjunto de dados que os agentes econômicos e o próprio banco central utilizam em suas decisões, o que pode levar a uma má especificação do choque monetário \citep[]{villaverde}. Como consequência, os SVARs tendem a subestimar a magnitude do impacto da política monetária e a gerar respostas contraintuitivas em algumas variáveis \citep[]{alessi}. Para superar essa limitação, este estudo emprega um Modelo de Fatores Dinâmicos Estruturais (SDFM), que extrai fatores comuns de um amplo conjunto de variáveis macroeconômicas e financeiras, aproximando o espaço de informação relevante e permitindo uma identificação mais robusta dos choques \citep[]{STOCK2016415, alessi}.

A identificação do choque monetário é realizada por meio de um instrumento externo construído a partir das surpresas de juros observadas em dias de reunião do Copom, combinando a ortogonalização por informação predeterminada de \cite{bauer2023} e a regra de sinal de \cite{jarocinski2020}, que separa choques de política de choques de informação do banco central. A relevância do instrumento é avaliada pela estatística de Wald de \cite{montielolea}, complementar ao F robusto de primeiro estágio.

O modelo é estimado utilizando um painel com 106 séries macroeconômicas e financeiras brasileiras, abrangendo atividade econômica, mercado de trabalho, inflação, crédito, agregados monetários e preços de ativos entre janeiro de 2013 e setembro de 2025. Esse amplo conjunto de informações permite identificar os choques monetários de forma mais robusta do que em VARs de pequena dimensão e acompanhar sua transmissão simultânea para diversas classes de ativos, câmbio, prêmio de risco soberano, curva de juros, crédito e mercado acionário, dimensão conjunta que desenhos de menor escala não permitem capturar.

Nossos resultados também diferem das evidências recentes para o Brasil obtidas por identificação via heterocedasticidade em frequência diária, que rejeitam a hipótese de dominância fiscal \cite{goncalves2025}. Essa divergência sugere que a importância do canal de prêmio de risco pode depender da frequência dos dados e da estratégia de identificação utilizada, uma questão que discutimos ao longo do artigo.

O restante deste artigo está organizado da seguinte forma. A Seção 2 revisa a literatura sobre identificação de choques monetários e canais de transmissão para preços de ativos. A Seção 3 descreve o modelo de fatores dinâmicos estruturais, a base de dados e a estratégia de identificação por instrumento externo. A Seção 4 apresenta os resultados, com ênfase no canal de câmbio e risco soberano. A Seção 5 reúne os diagnósticos de exogeneidade e invertibilidade e testa se o filtro de sinal seleciona notícias de risco soberano ou do FOMC. A Seção 6 conclui.

%A compreensão dos canais de transmissão da política monetária é um dos temas centrais da macroeconomia, pois as decisões do banco central afetam diretamente as condições de crédito, as expectativas dos agentes e a atividade econômica. Dentre os diversos canais, a resposta dos preços dos ativos financeiros — como ações, títulos e a taxa de câmbio — é particularmente relevante, pois reflete a reação imediata dos mercados e serve como um importante mecanismo de transmissão para a economia real. A literatura recente tem dedicado atenção considerável a quantificar essa relação, utilizando modelos cada vez mais sofisticados para isolar o impacto de choques monetários na precificação de risco (e.g., \cite{BEKAERT2013771}, \cite{alessi}). No contexto de economias emergentes como o Brasil, caracterizadas por maior volatilidade e particularidades institucionais, a análise desses efeitos torna-se ainda mais crucial para a condução da política monetária e a gestão de portfólios.

%A correta identificação desses efeitos depende crucialmente da informação disponível para o econometrista. Modelos tradicionais, como os Vetores Autorregressivos Estruturais (SVAR), frequentemente sofrem do problema da não-fundamentalidade, pois não conseguem incorporar o vasto conjunto de dados que os agentes econômicos e o próprio banco central utilizam em suas decisões, o que pode levar a uma má especificação do choque monetário \cite{villaverde}. Estudos mostram que, por consequência, os SVARs tendem a subestimar a magnitude do impacto da política monetária e a gerar respostas contraintuitivas em algumas variáveis \cite{alessi}.

%Para superar essa limitação, este estudo emprega um Modelo de Fatores Dinâmicos Estruturais (SDFM). Ao extrair fatores comuns de um grande conjunto de variáveis macroeconômicas e financeiras, o SDFM permite uma identificação mais robusta dos choques, capturando a informação relevante de maneira parcimoniosa \cite{STOCK2016415}. A principal contribuição deste trabalho é, portanto, fornecer uma quantificação mais precisa dos canais de transmissão da política monetária no Brasil, mostrando que os efeitos sobre os preços dos ativos são mais fortes e rápidos do que a literatura baseada em modelos de menor dimensão tipicamente encontra. O restante do artigo detalha a metodologia, os dados e aprofunda a análise dos resultados.

%O desenho é o de \cite{alessi}, aplicado a uma economia emergente com regime de dívida distinto do europeu e do americano. O instrumento é construído sobre surpresas de DI em dias de reunião do Copom. Ele combina a ortogonalização por informação predeterminada de \cite{bauer2023} e a regra de sinal de \cite{jarocinski2020}. A relevância é avaliada pela estatística de \cite{montielolea} e não pela F de primeiro estágio.

\section{Revisão de literatura}
\label{sec:literatura}

Dois problemas separam as estimativas do efeito da política monetária sobre preços de ativos. O primeiro é o tamanho do conjunto de informação do modelo. O segundo é a identificação do choque. A literatura tratou os dois em linhas quase independentes.

Um VAR de pequena dimensão raramente contém a informação que o banco central observa. \cite{villaverde} formalizam a condição sob a qual as inovações de um VAR geram o mesmo espaço que os choques estruturais. Ela falha em modelos de equilíbrio geral usuais. Quando falha, a reação sistemática do banco central a informação omitida aparece como choque de política.

A resposta da literatura foi ampliar o painel. \cite{bernanke} extraem fatores de um conjunto amplo de séries e os incluem no VAR, o modelo FAVAR. \cite{STOCK2016415} tratam o caso geral, em que o modelo de fatores dinâmicos substitui o VAR e a identificação opera sobre as inovações fatoriais. Trabalhos posteriores aplicam a mesma lógica à transmissão global da política americana \cite{miranda} e ao sentimento do investidor \cite{lutz}.

O argumento tem uma pré-história empírica. \cite{bagliano1998} confrontam os choques de um VAR monetário com medidas de surpresa extraídas de mercados futuros. Essas medidas usam o conjunto de informação de que os mercados dispõem, ao contrário da função de reação estimada. Ainda assim produzem funções de resposta que não diferem estatisticamente das do VAR de referência.

\cite{alessi} isolam o efeito do conjunto de informação mantendo a identificação fixa. Os autores estimam um VAR de quatro variáveis e um modelo de fatores dinâmicos de grande escala sobre os mesmos dados. O instrumento de alta frequência, a ordem de defasagens e a normalização de 50 pontos-base no yield soberano de 2 anos são os mesmos nos dois modelos. Os painéis reúnem 88 séries para a área do euro e 95 para os Estados Unidos.

As diferenças entre os dois modelos são grandes. Nos dados americanos, a resposta do prêmio excedente dos títulos corporativos no impacto é o dobro no modelo de fatores. Os preços de imóveis caem cerca de 1,5\% em dois anos no modelo de fatores e sobem no VAR. A maioria das especificações do VAR produz aumento da produção industrial após um aperto. Como apenas o conjunto de informação varia, a diferença não pode ser atribuída à hipótese de identificação.

Os modelos teóricos que motivam essa literatura tratam o preço dos ativos como canal de transmissão. \cite{castelnuovo} incorporam rigidez nominal de salários e heterogeneidade de investidores em um modelo DSGE, e obtêm efeito da riqueza financeira sobre o consumo. \cite{Cooley} mostram, em equilíbrio geral com firmas heterogêneas, que a estrutura financeira da firma condiciona sua resposta ao choque monetário. As previsões que o trabalho empírico procura medir são a reprecificação imediata e a resposta desigual entre classes de ativos.

A identificação por instrumento externo trata o segundo problema sem impor restrições de exclusão sobre a matriz de impacto. \cite{mertensravn2013} são a referência de origem do estimador para o caso geral com múltiplos instrumentos. Os autores tratam as séries narrativas de mudanças tributárias não como observações do choque, mas como variáveis correlacionadas com ele. Registros históricos raramente são inequívocos. Uma série tratada como o próprio choque produz estimativas atenuadas pelo erro de medida, e o mesmo erro não compromete sua validade como instrumento.

As condições que enunciam são relevância, $E[m_t \varepsilon_{1t}'] = \Phi$ não singular, e exogeneidade, $E[m_t \varepsilon_{2t}'] = 0$, sobre os demais choques estruturais. Os autores também medem a qualidade da proxy. As confiabilidades estimadas de suas duas séries narrativas são $0{,}30$ e $0{,}69$, o que corresponde a correlações de $0{,}55$ e $0{,}83$ com os choques latentes. A leitura que fazem é que o erro de medida em séries narrativas é grande na prática.

\cite{stockwatson2018} separam o que cada estimador exige. O SVAR-IV projeta o instrumento sobre as inovações do VAR e propaga a resposta pela dinâmica estimada. Exige relevância, exogeneidade contemporânea e invertibilidade, esta última a hipótese de que os choques estruturais podem ser recuperados de valores correntes e passados das observáveis.

O LP-IV estima uma regressão de variável instrumental por horizonte e dispensa a invertibilidade. Em troca, exige exogeneidade também em relação a choques passados e futuros. A parte restritiva é a que envolve choques passados.

A saída fácil desse dilema não existe. Incluir defasagens genéricas das observáveis como controles, para salvar um instrumento correlacionado com choques passados, equivale a supor invertibilidade. Nesse caso o SVAR-IV é o estimador mais eficiente. Os mesmos autores apontam o modelo de fatores como caminho para a invertibilidade, com a ressalva de que ampliar o número de variáveis não a garante.

\cite{gertler2015} levam o instrumento externo para a política monetária. A surpresa é medida sobre contratos futuros de juros, em janela de 30 minutos em torno dos anúncios do Federal Open Market Committee. Os autores separam instrumento de política, a taxa overnight, de indicador de política, uma taxa de prazo mais longo cuja inovação mensal incorpora a revisão das expectativas de juros.

A escolha do indicador é empírica. Os autores adotam a taxa de 1 ano, e não a de 2 anos que preferem conceitualmente. A razão é que nenhuma combinação de surpresas alcança primeiro estágio utilizável para a taxa de 2 anos. A F de primeiro estágio da taxa de 1 ano é $21{,}5$, contra o limiar convencional de 10. O resultado substantivo é que uma alta de 20 pontos-base nessa taxa eleva as taxas de crédito corporativo em cerca de 15 pontos-base. Quase todo o aumento vem de prêmio de prazo e spread de crédito.

Duas críticas posteriores mostram que a surpresa bruta não satisfaz a exogeneidade. \cite{jarocinski2020} observam que cerca de um terço dos anúncios do Federal Open Market Committee desde 1990 traz juros e ações subindo juntos. O co-movimento é incompatível com um choque de política pura. Os autores atribuem esses dias a um choque de informação do banco central. A separação usa o sinal do co-movimento e zera a surpresa nos dias de sinal positivo. Com ela, a resposta dos juros fica menos persistente e a queda do nível de preços mais pronunciada, o que atenua o \textit{price puzzle}.

\cite{bauer2023} chegam ao mesmo sintoma por outra via. Os autores mostram que a surpresa de alta frequência é previsível por notícias públicas anteriores ao anúncio, com $R^2$ entre $0{,}12$ e $0{,}20$. A causa que apontam é que o mercado subestima a resposta do Federal Reserve às notícias. Nenhuma informação privada do banco central é necessária para gerar o padrão.

O remédio é regredir a surpresa sobre essa informação predeterminada e usar o resíduo como instrumento. A correção vale para identificação em VAR e em \textit{local projections}, e é dispensável em estudos de evento com preços de ativos. As duas camadas atacam problemas distintos, contaminação por informação e previsibilidade, e podem ser aplicadas em conjunto, embora a literatura não tenha formalizado essa combinação..

\cite{montielolea} tratam a inferência quando o instrumento é fraco. A função de resposta identificada por instrumento externo é uma razão entre quantidades assintoticamente normais, e o denominador dessa razão é a própria relevância. Quando ela é pequena, o estimador de plug-in é inconsistente e viesado na direção da resposta que a decomposição de Cholesky produziria com o choque ordenado em primeiro lugar. A consequência prática é direta. A proximidade entre uma estimativa recursiva e uma por instrumento externo não corrobora a ordenação recursiva, a menos que a força do instrumento seja verificada.

Paralelamente à discussão sobre identificação, uma literatura recente documenta que a resposta cambial a choques de juros pode se afastar da previsão convencional da paridade descoberta quando fragilidades estruturais tornam o prêmio de risco endógeno ao próprio choque. \cite{dalgicozhan2026}, por exemplo, mostram que em economias com descasamento cambial nos balanços bancários e faturamento de exportações em dólar, um aperto monetário externo eleva o prêmio de risco cambial exigido pelos investidores, atenuando ou revertendo a apreciação que a paridade descoberta prescreveria. Embora a fonte de fragilidade nesse caso seja financeiro-cambial, e não fiscal, o resultado ilustra que a dominância do canal de prêmio de risco sobre o diferencial de juros não é uma singularidade brasileira, mas um fenômeno documentado em diferentes contextos institucionais. Se um mecanismo análogo opera no Brasil por meio do canal fiscal, contudo, é uma questão que a literatura recente ainda não resolve.

Para o Brasil, \cite{goncalves2025} identificam choques monetários por heterocedasticidade, no desenho de \cite{rigobon2003}. A comparação é entre semanas com e sem reunião do Copom. Os autores usam variações diárias de quarta para quinta-feira, em 768 observações entre 2009 e 2024. As surpresas vêm de contratos de DI em quatro vencimentos, e as expectativas de inflação de títulos indexados.

Um aperto de 100 pontos-base reduz a expectativa de inflação de 1 ano entre $0{,}20$ e $0{,}27$ ponto percentual, e a de 5 anos entre $0{,}53$ e $0{,}69$ ponto percentual. O real aprecia entre $3{,}4\%$ e $5{,}6\%$, e o CDS de 5 anos não responde. A leitura dos autores é que a política monetária ancora expectativas mesmo com dívida elevada, o que rejeita a hipótese de dominância fiscal em frequência diária. As estimativas mensais deste trabalho divergem no sinal do câmbio e na resposta do risco soberano, desacordo retomado na seção de resultados.


\section{Metodologia e Estratégia Empírica}
\label{sec:metodologia}

\subsection{O Modelo de Fatores Dinâmicos Estruturais}

O modelo empírico é um Modelo de Fatores Dinâmicos Estruturais (SDFM), representado por três equações:
\begin{align}
    \overset{n\times 1}{X_t} &= \overset{n\times r}{\Lambda} \overset{r\times 1}{F_t} + \overset{n\times 1}{e_t}\label{eq:4}\\[4pt]
    \overset{r\times r}{\Phi(L)}\overset{r\times 1}{F_t} &= \overset{r\times q}{G} \overset{q\times 1}{\eta_t} \label{eq:5}\\[4pt]
    \overset{q\times 1}{\eta_t} &= \overset{q\times q}{H} \overset{q\times 1}{\varepsilon_t} \label{eq:6}
\end{align}
A \ref{eq:4} é a equação de observação: o vetor de $n$ variáveis observáveis $X_t$ decompõe-se em um componente comum, dado por $r$ fatores latentes $F_t$ e pela matriz de cargas $\Lambda$, e um componente idiossincrático $e_t$. A \ref{eq:5} descreve a dinâmica dos fatores, com $\Phi(L) = I - \Phi_1 L - \dots - \Phi_p L^p$ um polinômio de defasagens de ordem $p$ e $\eta_t$ um vetor de $q$ inovações fatoriais. A \ref{eq:6} liga essas inovações aos $q$ choques estruturais $\varepsilon_t$ pela matriz $H$. A ordem do VAR dos fatores é fixada em $p = 6$, suficiente para acomodar a persistência das séries mensais brasileiras sem sobreparametrizar a dinâmica.

Combinando as três equações, o modelo assume a forma de média móvel que escreve as observáveis em termos dos choques estruturais:
\begin{align}\label{eq:7}
    X_t=\Lambda \Phi(L)^{-1}GH\varepsilon_t + e_t
\end{align}
A estimação requer determinar o número de fatores estáticos $r$ e dinâmicos $q$ e normalizar as matrizes $\Lambda$, $G$ e $H$.

O conjunto de séries é não-estacionário por construção, pois reúne séries em nível, séries em taxa e séries já diferenciadas, com ordens de integração distintas. Extrair fatores diretamente sobre esse conjunto faria com que as séries de maior variância dominassem a decomposição. Seguindo \cite{barigozzi2016non}, cada série sem tendência $X_{it}$ é dividida pelo desvio padrão de sua primeira diferença $s_{iy}$:
\begin{equation}
    Z_{it} = \frac{X_{it}}{s_{iy}}
\end{equation}
A padronização equaliza a escala das flutuações sem diferenciar as séries, o que preserva a informação de baixa frequência usada pelo modelo. Por essa razão os critérios de \cite{bai-ng} não entram na forma padrão: eles pressupõem séries estacionárias. O tratamento de \cite{barigozzi2016non} mantém a estrutura de fatores válida sob não-estacionariedade.

O uso de um grande conjunto de séries responde ao problema de não-fundamentalidade dos choques estruturais \cite{villaverde}. Um VAR de pequena dimensão raramente contém a informação que os agentes e o banco central de fato observam, e o choque recuperado nesse espaço restrito é uma combinação de choques estruturais, não o choque de política monetária isolado. Ao extrair fatores comuns de um conjunto amplo de variáveis, o SDFM aproxima o espaço de informação relevante e recupera choques fundamentais, como mostram \cite{alessi} ao comparar o SDFM a um VAR de referência sob o mesmo instrumento.

\subsection{Base de Dados}

A base reúne 111 séries mensais para a economia brasileira entre janeiro de 2013 e setembro de 2025, o que corresponde a 153 observações alinhadas e 147 inovações fatoriais após as seis defasagens do VAR. O início em 2013 acompanha a mudança metodológica da PNAD Contínua no mercado de trabalho. As fontes são oficiais (Banco Central do Brasil, IBGE, B3, Tesouro Nacional e ANBIMA), somadas a indicadores de risco soberano e mercado externo (EMBI, CDS de cinco anos e MSCI) e a um índice de incerteza de política econômica (EPU). Os blocos cobrem atividade real, mercado de trabalho, inflação, commodities, agregados monetários e de crédito, e preços de ativos. Como a interpretação dos resultados envolve risco soberano, o painel inclui um bloco fiscal, com a dívida bruta do governo geral, a dívida líquida do setor público e o resultado primário consolidado, e um bloco de expectativas, com as medianas da pesquisa Focus para o IPCA acumulado em doze meses e para a Selic, o PIB e o câmbio no fim do ano seguinte.

As taxas de juros de maturidade fixa (3 meses, 6 meses, 1, 2, 5 e 10 anos) vêm do ajuste do modelo de \cite{svensson1994estimating} sobre os contratos de DI futuro em cada data. As séries nominais entram em logaritmo e, quando os testes indicam componente sazonal, são ajustadas pelo procedimento X-13ARIMA-SEATS \cite{sax2018seasonal}. A lista completa de séries, fontes e transformações está no Anexo \ref{anexo-A} (\autoref{tab:lista_variaveis}).

\subsection{Identificação por instrumento externo}

A identificação do choque de política monetária usa um instrumento externo, no arcabouço de proxy-SVAR de \cite{stockwatson2018} e \cite{gertler2015}, transposto para o DFM como em \cite{alessi}. A abordagem dispensa a ordenação recursiva entre os fatores. No lugar dela entra uma variável observada fora do modelo, correlacionada com o choque de interesse e ortogonal aos demais choques estruturais.

Seja $z_t$ o instrumento e $\varepsilon_t^{p}$ o choque de política monetária. O instrumento é válido sob duas condições: relevância, $E[z_t \varepsilon_t^{p}] \neq 0$, e exogeneidade, $E[z_t \varepsilon_t^{j}] = 0$ para todo choque estrutural $j \neq p$ \cite{stockwatson2018}. Essas duas são exatamente a condição que \cite{stockwatson2018} associam ao estimador SVAR-IV, e ambas são contemporâneas. A condição de exogeneidade em relação a choques passados e futuros, que os mesmos autores chamam de lead-lag, é mais forte e pertence à alternativa por projeções locais; ela não é exigida aqui porque as inovações fatoriais já são ortogonais ao passado dos fatores e não contêm choques futuros. O preço dessa dispensa é a invertibilidade, discutida na Seção \ref{sec:literatura} e diagnosticada na Seção \ref{sec:invertibilidade}. Sob essas condições, a projeção das inovações fatoriais $\eta_t$ sobre o instrumento identifica, a menos de escala, a coluna de $H$ associada ao choque de política monetária. Empilhando $\eta_t$ e $z_t$ na amostra alinhada, essa coluna é estimada por $\hat{H}_{\cdot p} = (Z'\eta)/(Z'Z)$. O alinhamento temporal segue a implementação computacional de \cite{alessi}.

O choque é normalizado para elevar o yield de 6 meses em 50 pontos-base no impacto, ou $0{,}005$ em proporção decimal. A escolha do yield de 6 meses como variável de política, em vez da Selic overnight, segue a distinção de \cite{gertler2015} entre o instrumento de política, a taxa overnight, e o indicador de política, uma taxa de prazo mais longo que incorpora a revisão das expectativas de juros embutida na decisão. No caso brasileiro a distinção é mais aguda. As decisões do Copom saem após o fechamento do mercado, e a surpresa é medida na curva de DI entre a quarta e a quinta-feira, não na taxa overnight. A Selic mensal, uma média do overnight dentro do mês, quase não embute essa surpresa. Sob o diagnóstico de relevância de \cite{montielolea}, a normalização pela Selic é a mais fraca das duas em toda a varredura de especificações da amostra completa, onde a Wald $\xi_{mp}$ máxima é $5{,}20$ contra $10{,}73$ da normalização pelo yield de 6 meses. Na janela pré-COVID as duas sobem, para $21{,}54$ e $22{,}24$, mas essa janela é marginalmente instável e por isso não decide a comparação.

\subsection{Construção do instrumento}

O instrumento primário combina três camadas: a surpresa de juros em dias de reunião do Copom, medida como em \cite{gertler2015}, a ortogonalização em relação à informação pública anterior ao anúncio, de \cite{bauer2023}, e o filtro de sinal de \cite{jarocinski2020}, que separa choques de política de choques de informação.

A primeira camada mede a surpresa como a variação do DI futuro no vértice de 126 dias úteis, cerca de 6 meses, entre o fechamento da quarta-feira da reunião e o da quinta-feira seguinte. A decisão do Copom é anunciada após o fechamento do mercado, de modo que a janela de um dia captura a primeira sessão em que os preços incorporam a decisão. A amostra tem 95 reuniões entre 2013 e 2025. As surpresas são somadas dentro de cada mês, o esquema de agregação de \cite{jarocinski2020}, e o instrumento vale zero nos meses sem reunião.

A segunda camada remove o componente previsível da surpresa. \cite{bauer2023} mostram que surpresas de alta frequência são correlacionadas com informação pública anterior ao anúncio, porque o mercado subestima a resposta do banco central às notícias. Uma surpresa previsível viola a exogeneidade do instrumento. Seguindo o desenho dos autores, a surpresa é regredida sobre preditores medidos até o fechamento da quarta-feira: uma tendência temporal, as variações acumuladas em 65 pregões do Ibovespa, do S\&P 500, do VIX, do petróleo Brent, do câmbio BRL/USD e da inclinação da curva de DI (2 anos menos 3 meses), e as revisões em 20 pregões das medianas da pesquisa Focus para o IPCA acumulado em 12 meses e para a Selic no fim do ano seguinte. A inclusão do retorno do Ibovespa entre os preditores também tem motivo de identificação, porque a regra de política estimada por \cite{castelnuovo} responde a condições do mercado acionário, canal que a ortogonalização precisa remover da surpresa antes do filtro de sinal. O resíduo dessa regressão é a surpresa ortogonalizada. O $R^2$ é $0{,}024$, na faixa esperada. \cite{bauer2023} obtêm entre $0{,}12$ e $0{,}20$ incluindo surpresas de divulgações macroeconômicas que não existem em frequência diária para o Brasil, e um $R^2$ alto indicaria vazamento de informação contemporânea na regressão. A exigência de que todos os preditores sejam predeterminados disciplina a especificação. Variáveis domésticas contemporâneas à janela do anúncio respondem ao próprio choque e enviesariam o resíduo. Variáveis predeterminadas, domésticas ou não, são informação disponível antes da decisão, exatamente o que a correção de \cite{bauer2023} usa.

A terceira camada separa o choque de política monetária do choque de informação do banco central. \cite{jarocinski2020} distinguem os dois pelo sinal do co-movimento entre a surpresa de juros e a do mercado de ações no mesmo dia. Juros em alta com ações em queda indicam um choque de política. Juros e ações em alta juntos indicam que o comunicado revelou informação favorável sobre a economia. O filtro classifica cada reunião pelo co-movimento dos resíduos ortogonalizados da surpresa de DI e do retorno do Ibovespa, mantém os dias de co-movimento negativo e zera os demais. Das 95 reuniões, 62 são classificadas como choque de política e 33, ou $34{,}7\%$, como informação, proporção próxima do terço que os autores encontram para o Fed. Como os preditores da segunda camada são predeterminados, a classificação preserva o co-movimento observado no dia do anúncio, em vez de depender de resíduos construídos com informação da própria janela. Essa é a variante fiel às duas metodologias de origem, a ortogonalização por preditores predeterminados de \cite{bauer2023} e a regra de sinal com agregação mensal por soma de \cite{jarocinski2020}. A principal diferença em relação ao filtro literal de \cite{jarocinski2020}, que classifica e agrega a surpresa bruta, é o uso dos resíduos predeterminados no lugar dela, sem efeito prático sobre a classificação pela razão acima. 

\subsection{Seleção de $(r, q)$ e estimação}

A especificação de produção é $(r, q) = (5, 5)$ com $p = 6$. O número de fatores estáticos vem do critério de informação IC2 de \cite{bai-ng} na versão padronizada de \cite{barigozzi2016non}, que é a variante admissível sob a não-estacionariedade do conjunto e seleciona $r = 5$ tanto na amostra completa quanto na janela pré-COVID. O número de fatores dinâmicos é fixado em $q = 5$ como escolha operacional associada a esse $r$, e não por um critério próprio, de modo que a fundamentação de $q$ permanece aberta. Nenhuma das duas dimensões foi escolhida pela célula de maior força do instrumento, porque selecionar $(r, q)$ pela estatística de relevância converteria um diagnóstico em critério de especificação, procedimento que \cite{montielolea} desaconselham.

Os fatores estáticos são extraídos por componentes principais sobre o conjunto de séries padronizado e a dinâmica segue o VAR(6). A estimativa pontual das funções de impulso-resposta usa os coeficientes de mínimos quadrados ordinários, fiel à implementação de \cite{alessi}. A incerteza vem de um \textit{bootstrap wild} \cite{goncalveskilian2004}, com 800 replicações, robusto à heterocedasticidade dos resíduos. A correção de viés de pequena amostra de \cite{kilian1998small} entra apenas no processo gerador do bootstrap, não na estimativa pontual. As funções de impulso-resposta cobrem horizontes de 0 a 48 meses, com bandas de 68\% e 90\%.

\subsection{Relevância do instrumento}

A validade da inferência no proxy-SVAR depende da força do instrumento \cite{montielolea}. Com instrumento fraco, o estimador de impacto é uma razão de normais correlacionadas, com viés na direção da identificação recursiva e cobertura incorreta das bandas, de modo que tanto a magnitude quanto a largura do intervalo perdem a interpretação usual.

Os dois diagnósticos de relevância propostos por \cite{montielolea} são calculados na mesma direção em que o choque é normalizado. A Wald $\xi_{mp}$ projeta o primeiro estágio na combinação das inovações fatoriais que determina o impacto sobre o yield de 6 meses, com erros padrão de Eicker-White e com o instrumento residualizado nas defasagens do VAR dos fatores. O $F^{rob}_{mp}$ é o quadrado da estatística $t$ HC1 do instrumento na regressão dessa mesma combinação sobre o instrumento e os mesmos controles, de modo que as duas estatísticas diferem pela forma de estimar a variância e não pelo objeto que projetam.

Na especificação de produção, $\xi_{mp}$ vale $6{,}27$ na amostra completa e $10{,}99$ na janela pré-COVID, enquanto $F^{rob}_{mp}$ vale $10{,}12$ e $9{,}75$ (\autoref{tab:first_stage}). As duas janelas dão evidência mista, e por motivos invertidos, uma vez que na amostra completa é o $F$ robusto que passa de 10 com a Wald abaixo, enquanto na pré-COVID é a Wald que passa com o $F$ abaixo. Essa segunda configuração deve algo à severidade da correção HC1 com 78 observações efetivas e 30 defasagens fatoriais como controles. O valor 10 é uma referência convencional para o $F$ de primeiro estágio homoscedástico do 2SLS, e não um valor crítico que \cite{montielolea} derivem para a Wald robusta a heterocedasticidade usada aqui, então ele orienta a leitura sem constituir um teste. A janela pré-COVID tem ainda raiz máxima da companion igual a $1{,}000202$, o que a torna diagnóstico de força e estabilidade, não fonte de leitura dinâmica. Como \cite{montielolea} desaconselham selecionar o procedimento inferencial por um pré-teste de força, as bandas do \textit{wild bootstrap} são lidas com cautela nas duas janelas.

A \autoref{tab:first_stage} decompõe esses números pelas três camadas da construção do instrumento, e a força não cresce monotonicamente com elas. A ortogonalização em relação à informação predeterminada, tomada isoladamente, reduz $\xi_{mp}$ nas duas janelas, de $5{,}27$ para $4{,}37$ na amostra completa e de $14{,}86$ para $13{,}98$ na pré-COVID. O filtro de sinal move as duas janelas em direções opostas. Na amostra completa ele eleva $\xi_{mp}$ de $4{,}37$ para $6{,}27$ e $F^{rob}_{mp}$ de $4{,}32$ para $10{,}12$, ainda que reduza de 90 para 60 os meses com instrumento não-nulo. Na janela pré-COVID ele derruba $\xi_{mp}$ de $13{,}98$ para $10{,}99$, com $F^{rob}_{mp}$ quase parado, em $9{,}75$ contra $9{,}10$, de modo que nessa janela a surpresa bruta é a mais forte das três variantes pela Wald. A variante de produção não é, portanto, a de maior força em toda janela. Ela é a que reproduz as duas metodologias de origem, a ortogonalização por preditores predeterminados de \cite{bauer2023} e a regra de sinal de \cite{jarocinski2020}, e escolhê-la pela estatística de relevância seria o mesmo pré-teste que \cite{montielolea} desaconselham e que a subseção anterior já evita na seleção de $(r, q)$.

% tab:first_stage — números transcritos de dois artefatos gerados em 2026-08-13,
% ambos na produção de 111 séries com (r,q) = (5,5):
%   amostra completa (n = 147): output/instrument/instrument_diagnostics_report.md, secao 1
%   janela pre-COVID (n = 78):  output/instrument/mosw_strength_grid.csv, linhas sample=pre_covid, r=5, q=5
% O CSV da grade nao carrega beta/EP/p, e por isso o bloco de primeiro estagio
% existe so para a amostra completa. Reconferir contra esses dois arquivos antes
% de qualquer republicacao.
\begin{table}[H]
\begin{center}
\caption{\textbf{Relevância do instrumento por camada de construção}\\{Cada linha acrescenta uma camada à construção do instrumento, na ordem descrita na subseção Construção do instrumento: a surpresa do DI no vértice de 126 dias úteis na janela de um dia em torno da reunião do Copom, a ortogonalização em relação à informação pública predeterminada e o filtro de sinal, que zera as reuniões classificadas como choque de informação. A última linha é a variante de produção. A amostra completa vai de 2013:01 a 2025:09 e a janela pré-COVID de 2013:01 a 2019:12, com $147$ e $78$ inovações fatoriais depois das $p = 6$ defasagens. \emph{Meses $z \neq 0$} conta as observações em que o instrumento agregado é não-nulo. $\hat{\beta}$ e o erro padrão de Eicker-White HC1 vêm da regressão de primeiro estágio de $c_{mp}'\eta_t$ sobre o instrumento, com as defasagens do VAR dos fatores como controles, e estão multiplicados por $10^{4}$; $p$ é o valor-$p$ bicaudal correspondente. $\xi_{mp}$ e $F^{rob}_{mp}$ são as duas estatísticas definidas no corpo desta subseção, calculadas na mesma direção de normalização. O bloco de primeiro estágio é reportado apenas na amostra completa. O impacto estimado sobre o yield de 6 meses é positivo nas três variantes.}}
\label{tab:first_stage}
\footnotesize
\setlength{\tabcolsep}{5pt}
\begin{tabular}{@{}l c c c c c c c c@{}}
\toprule
 & \multicolumn{6}{c}{Amostra completa ($n = 147$)} & \multicolumn{2}{c}{Pré-COVID ($n = 78$)} \\
\cmidrule(lr){2-7} \cmidrule(lr){8-9}
Variante & Meses $z \neq 0$ & $\hat{\beta} \times 10^{4}$ & EP\,$\times\,10^{4}$ & $p$ & $\xi_{mp}$ & $F^{rob}_{mp}$ & $\xi_{mp}$ & $F^{rob}_{mp}$ \\
\midrule
Surpresa bruta do DI                    & 87 & $0{,}681$ & $0{,}287$ & $0{,}020$ & $5{,}27$ & $5{,}61$  & $14{,}86$ & $9{,}71$ \\
\quad + ortogonalização predeterminada  & 90 & $0{,}612$ & $0{,}295$ & $0{,}040$ & $4{,}37$ & $4{,}32$  & $13{,}98$ & $9{,}10$ \\
\quad + filtro de sinal (produção)      & 60 & $1{,}150$ & $0{,}361$ & $0{,}002$ & $6{,}27$ & $10{,}12$ & $10{,}99$ & $9{,}75$ \\
\bottomrule
\end{tabular}
\end{center}
\end{table}



\section{Resultados}
\label{sec:resultados}

Um aperto de 50 pontos-base no DI de 6 meses eleva o DI de 5 anos em $77{,}6$ pontos-base, deprecia o BRL/USD em $3{,}84\%$ da média amostral e aumenta o EMBI+ e o CDS de 5 anos em $26{,}2$ e $32{,}5$ pontos-base no impacto. O movimento conjunto da curva, do câmbio e do risco soberano é o principal resultado da estimação com 111 séries e $(r,q,p)=(5,5,6)$. Como a estatística de força no espaço dos fatores é $\xi_{mp}=6{,}27$ na amostra completa, as bandas de 68\% e 90\% descrevem a incerteza das respostas, mas não sustentam classificações binárias dos resultados.

\subsection{Estrutura a termo}

A curva inteira sobe no impacto, e o repasse aumenta com a maturidade até 5 anos. Os vértices de 3 meses, 1 ano, 2 anos, 5 anos e 10 anos aumentam $38{,}7$, $63{,}3$, $74{,}3$, $77{,}6$ e $70{,}3$ pontos-base, respectivamente, enquanto o DI de 6 meses sobe os 50 pontos-base impostos pela normalização (\autoref{fig:curva}). A Selic mensal também responde, com alta de $23{,}8$ pontos-base no impacto, de modo que a resposta inicial abrange tanto a taxa \textit{overnight} quanto os contratos de maior maturidade.

As respostas alcançam valores maiores nos primeiros meses e depois revertem, mas o horizonte e a magnitude dessa reversão variam com a defasagem do VAR dos fatores. A reversão é, portanto, uma propriedade da dinâmica conjunta estimada e não uma confirmação independente de um canal de prêmio de prazo ou de risco.

\begin{figure}[H]
\begin{center}
    \includegraphics[scale=0.7]{fig_curva.pdf}
\caption{\textbf{Estrutura a termo}\\{Resposta a um choque contracionista normalizado em 50 pontos-base no DI de 6 meses. A linha preta mostra a estimativa pontual em pontos-base, e as áreas escura e clara mostram as bandas de 68\% e 90\% do \textit{wild bootstrap} com 800 réplicas. A banda do DI de 6 meses é degenerada no impacto porque cada réplica recebe a mesma normalização. A Selic é a taxa \textit{overnight} mensal.}}
%\footnotesize{Source: Own elaboration, based on data from Central Bank.}
\label{fig:curva}
\end{center} 
\end{figure}

\subsection{Câmbio e risco soberano}\label{sec:cambio}
O real se deprecia no impacto contra as cinco moedas observadas. No BRL/USD, o aumento de R\$ $0{,}1579$ por dólar equivale a $3{,}84\%$ quando dividido pela média amostral de R\$ $4{,}1117$ por dólar. As demais respostas, calculadas da mesma forma, são $3{,}00\%$ no BRL/EUR, $3{,}67\%$ no BRL/CNY, $3{,}28\%$ no BRL/INR e $4{,}48\%$ no BRL/ARS. Essa conversão depende do nível usado como denominador, mas a concordância entre as cinco taxas indica que o resultado pertence ao real e não apenas ao par com o dólar.

A depreciação ocorre junto com aumentos de $26{,}2$ pontos-base no EMBI+ e $32{,}5$ pontos-base no CDS de 5 anos (\autoref{fig:cambio}). O co-movimento é compatível com transmissão por prêmio de risco, mas o desenho não decompõe o diferencial de juros nem identifica uma origem fiscal para esse prêmio. Por 100 pontos-base, a resposta cambial corresponde a aproximadamente $7{,}68\%$ de depreciação sob reescala linear, enquanto \cite{goncalves2025} estimam de $3{,}4\%$ a $5{,}6\%$ de apreciação. Essa normalização aritmética coloca o tamanho do choque na mesma base, mas não torna equivalentes estimativas com frequências, horizontes, estratégias de identificação e denominadores cambiais distintos. A resposta conjunta das duas medidas soberanas é consistente com o vínculo entre política monetária e precificação de risco documentado por \cite{BEKAERT2013771}.

\begin{figure}[H]
\begin{center}
    \includegraphics[scale=0.7]{fig_cambio_risco.pdf}
\caption{\textbf{Câmbio e risco soberano}\\{A linha preta mostra a estimativa pontual, e as áreas escura e clara mostram as bandas de 68\% e 90\%. As taxas de câmbio em nível são divididas pela respectiva média amostral e expressas em percentual, com valores positivos indicando depreciação do real. EMBI+ e CDS são expressos em pontos-base. Os dois últimos painéis são índices de commodities do Banco Central denominados em reais.}}
%\footnotesize{Source: Own elaboration, based on data from Central Bank.}
\label{fig:cambio}
\end{center} 
\end{figure}

\subsection{Atividade e mercado de trabalho}

A atividade recua no impacto, mas a magnitude varia entre produção, consumo, serviços e utilização de insumos. O IBC-Br cai $0{,}46\%$, a indústria de transformação cai $1{,}37\%$, os bens duráveis caem $4{,}14\%$ e os bens de capital caem $2{,}03\%$. No consumo, as vendas no varejo recuam $0{,}81\%$, enquanto as vendas de serviços diminuem $0{,}37\%$. A utilização da capacidade cai $0{,}12$ ponto percentual e as horas trabalhadas na indústria caem $0{,}85\%$ (\autoref{fig:atividade}).

A maior contração dos bens duráveis e a resposta menor dos serviços mostram que o efeito inicial não é uniforme entre setores. Nos meses seguintes, indústria e duráveis aprofundam a queda, enquanto serviços e capacidade apresentam trajetórias menores e menos regulares, de modo que o bloco informa heterogeneidade da transmissão, não uma cronologia única para toda a atividade.

\begin{figure}[H]
\begin{center}
    \includegraphics[scale=0.7]{fig_atividade.pdf}
\caption{\textbf{Atividade e horas trabalhadas}\\{A linha preta mostra a estimativa pontual, e as áreas escura e clara mostram as bandas de 68\% e 90\%. Índices de atividade, produção, vendas e horas são divididos pela respectiva média amostral e expressos em percentual. A utilização da capacidade instalada é expressa em pontos percentuais.}}
%\footnotesize{Source: Own elaboration, based on data from Central Bank.}
\label{fig:atividade}
\end{center} 
\end{figure}

\subsection{Crédito}

O crédito cresce de forma generalizada no impacto. O saldo total aumenta $0{,}57\%$, e as altas setoriais são $1{,}28\%$ no agropecuário, $1{,}42\%$ no transporte, $0{,}82\%$ na indústria, $0{,}79\%$ no comércio e $0{,}65\%$ na construção. O crédito à pessoa física aumenta $0{,}29\%$, a menor resposta inicial do bloco (\autoref{fig:credito}).

Essa expansão diminui ao longo do primeiro ano e se transforma em contração nas estimativas posteriores, em geral depois de dois anos. O padrão pode resultar de saques iniciais de linhas já contratadas seguidos pelo ajuste de balanços, mecanismo relacionado à heterogeneidade financeira estudada por \cite{Cooley}. As respostas, porém, não separam esse mecanismo de mudanças na composição do crédito ou de outras defasagens da atividade, e por isso não identificam um acelerador financeiro.

\begin{figure}[H]
\begin{center}
    \includegraphics[scale=0.7]{fig_credito.pdf}
\caption{\textbf{Crédito}\\{A linha preta mostra a estimativa pontual, e as áreas escura e clara mostram as bandas de 68\% e 90\%. Os saldos entram em log-nível, de modo que as respostas são log-pontos multiplicados por 100 e aproximam variações percentuais.}}
%\footnotesize{Source: Own elaboration, based on data from Central Bank.}
\label{fig:credito}
\end{center} 
\end{figure}

\subsection{Preços}

O IPP aumenta $0{,}468$ ponto percentual no impacto, enquanto o IGP-M sobe $0{,}128$. Os preços ao consumidor não apresentam a mesma alta imediata, pois o IPCA cheio cai $0{,}062$ ponto percentual, os núcleos EX0 e de médias aparadas ficam próximos de zero, o núcleo EX1 cai $0{,}028$ e o INPC cai $0{,}077$ (\autoref{fig:precos}). Essa diferença separa o aumento contemporâneo dos preços ao produtor da resposta posterior dos preços ao consumidor.

No sexto mês, o IPCA aumenta $0{,}089$ ponto percentual, e os núcleos EX0, de médias aparadas e EX1 sobem $0{,}090$, $0{,}068$ e $0{,}050$, respectivamente. O INPC sobe $0{,}059$ ponto percentual no mesmo horizonte, de modo que a resposta posterior é positiva em várias medidas, mas difere em magnitude e duração. A sequência é compatível com repasse cambial de custos, embora não identifique esse canal separadamente.

O instrumento usa o filtro de sinal de \cite{jarocinski2020}, mas a comparação entre as oito variantes do instrumento ainda não foi concluída. Por esse motivo, a seção se limita a descrever a sequência entre preços ao produtor e ao consumidor.

\begin{figure}[H]
\begin{center}
    \includegraphics[scale=0.7]{fig_precos.pdf}
\caption{\textbf{Preços}\\{A linha preta mostra a estimativa pontual, e as áreas escura e clara mostram as bandas de 68\% e 90\%. As séries de inflação são taxas mensais e as respostas são expressas em pontos percentuais da taxa. O índice de difusão é expresso em pontos percentuais da proporção de itens com aumento de preço.}}
%\footnotesize{Source: Own elaboration, based on data from Central Bank.}
\label{fig:precos}
\end{center} 
\end{figure}

\subsection{Ações}

Os sete índices da B3 têm retorno mensal negativo no impacto. As respostas são $-1{,}72\%$ no Ibovespa, $-2{,}61\%$ no SMLL, $-2{,}04\%$ no IDIV, $-2{,}56\%$ no IMOB, $-2{,}70\%$ no IFNC, $-0{,}07\%$ no IMAT e $-1{,}31\%$ no IFIX (\autoref{fig:acoes}). A faixa entre $-2{,}70\%$ e $-0{,}07\%$ mostra uma queda disseminada, mas também uma diferença ampla entre os índices.

O IMAT é a exceção econômica dentro desse movimento porque reúne empresas com receita denominada em dólar. A depreciação do real na mesma janela aumenta o valor dessa receita em moeda doméstica e é compatível com o amortecimento da queda, sem exigir uma resposta rápida dos preços de exportação. Nos horizontes posteriores, a figura continua mostrando retornos mensais, e não uma trajetória acumulada do nível dos índices.

\begin{figure}[H]
\begin{center}
    \includegraphics[scale=0.7]{fig_acoes.pdf}
\caption{\textbf{Índices de ações}\\{A linha preta mostra a resposta do retorno mensal em pontos percentuais, sem acumulação ao longo dos horizontes. As áreas escura e clara mostram as bandas de 68\% e 90\%.}}
%\footnotesize{Source: Own elaboration, based on data from Central Bank.}
\label{fig:acoes}
\end{center} 
\end{figure}

\section{Robustez}
\label{sec:robustez}

\subsection{Exogeneidade do instrumento}
\label{sec:exogeneidade}

A ortogonalização por informação opera no dia do anúncio e não garante que o instrumento, uma vez agregado em frequência mensal, seja imprevisível. O primeiro diagnóstico testa a previsibilidade mensal do instrumento em cinco especificações, todas representadas por
\begin{equation}
z_t = \alpha + \sum_{\ell=1}^{L} \boldsymbol{\beta}_\ell' \mathbf W_{t-\ell} + u_t,
\qquad
H_0\colon \boldsymbol{\beta}_1 = \cdots = \boldsymbol{\beta}_L = \mathbf 0,
\label{eq:predictability}
\end{equation}
em que $\mathbf W_t$ é o vetor-coluna de regressores e $\boldsymbol{\beta}_\ell$ reúne os coeficientes associados à defasagem $\ell$. Duas especificações usam os retornos das séries externas, com $L=3$ e $L=6$\footnote{O bloco externo contém a variação do índice VIX, o retorno do MSCI de mercados emergentes, a variação do índice de commodities metálicas do Banco Central e a variação do índice de incerteza de política econômica dos Estados Unidos, enquanto o câmbio USD/BRL entra como quinto regressor, de natureza doméstica.}. As outras três usam os cinco fatores estimados, com $L=1$, $L=3$ e $L=6$. Em cada especificação, a estatística de Wald robusta a heterocedasticidade testa conjuntamente a hipótese de que todos os vetores de coeficientes são nulos. Como o instrumento vale zero em 94 dos 156 meses do calendário, o valor-$p$ assintótico não é confiável, e os valores reportados vêm de \textit{wild bootstrap} sob a nula com 2.000 réplicas. Nenhuma das cinco especificações rejeita a imprevisibilidade, com $R^2$ entre $0{,}044$ e $0{,}263$ e valores-$p$ de \textit{bootstrap} entre $0{,}199$ e $0{,}749$. A escolha do \textit{bootstrap} decide o veredito: na especificação com os cinco fatores e $L=6$ o valor-$p$ assintótico é $0{,}007$ e o de \textit{bootstrap} é $0{,}400$, de modo que a inferência assintótica acusaria uma violação que o desenho apropriado a um instrumento censurado em zero não encontra. Duas especificações adicionais substituem os fatores em nível pelas inovações fatoriais defasadas, com valores-$p$ de $0{,}631$ em $L=6$ e $0{,}926$ em $L=3$.

Essa bateria testa a metade retrospectiva da condição de exogeneidade. Ela é a metade restritiva em \cite{stockwatson2018} e a que motiva a segunda camada de construção, porque é exatamente a previsibilidade documentada por \cite{bauer2023} que a ortogonalização remove no dia do anúncio, sem garantia de sobreviver à agregação mensal. Convém registrar o alcance do teste: como observam \cite{mertensravn2013}, a previsibilidade do instrumento não invalida por si só a identificação do proxy-SVAR, cujas condições são contemporâneas, de modo que o exercício mede a qualidade do instrumento e não a consistência do estimador. A única ressalva que não desaparece está no câmbio, cujas defasagens produzem valor-$p$ de $0{,}064$ em retornos com $L=6$, o único regressor do bloco externo abaixo de $0{,}10$, e que não figura entre os preditores predeterminados da segunda camada.

O segundo diagnóstico testa a denominação dos índices de commodities e é distinto das regressões de previsibilidade. Nos dois últimos painéis da \autoref{fig:cambio}, o índice de metais do Banco Central sobe $2{,}46\%$ da sua média amostral no impacto e exclui zero a 90\% até $h = 4$, embora uma economia pequena e aberta não mova o preço mundial dos metais. A denominação do índice permite testar se essa resposta vem de um componente global ou do câmbio.

O índice do Banco Central é medido em reais. Para cada índice de commodities $i$, a conversão
\[
C^{USD}_{it} = \frac{C^{BRL}_{it}}{S_t},
\]
em que $S_t$ é a taxa de câmbio em reais por dólar, separa o preço internacional da variação cambial. Em um painel aumentado com os três índices convertidos para dólares, estimado com 200 réplicas, o índice de metais em reais responde $+2{,}48\%$ no impacto e exclui zero nos cinco primeiros horizontes, enquanto o mesmo índice em dólares responde $-0{,}91\%$ e contém zero nos 25 horizontes testados. Os índices de agropecuária e de energia têm o mesmo padrão, e a correlação entre a variação do índice de metais e a do câmbio é $+0{,}450$ em reais contra $-0{,}409$ em dólares.

Um fator global moveria as duas denominações, mas a versão em dólares contém zero em todos os horizontes. O movimento do índice em reais herda, portanto, parte da resposta cambial, que é $+3{,}46\%$ no mesmo painel aumentado. Como a inclusão das três séries reduz $\xi_{mp}$ de $6{,}27$ para $5{,}77$, a comparação identifica o papel da denominação sem fornecer uma estimativa alternativa da magnitude de produção.

As três séries externas usadas como placebo passam nos dois níveis de leitura adotados na Seção 4. Na banda de 90\% nenhuma das três exclui zero em qualquer um dos 49 horizontes, e na banda de 68\% o único par que exclui zero é o do VIX em $h = 8$, somando 1 dos 147 pares das três séries.

\begin{figure}[H]
\begin{center}
    \includegraphics[scale=0.7]{fig_placebos.pdf}
\caption{\textbf{Placebos}\\{Séries externas que um choque de política monetária brasileiro não deve mover. Nenhum dos 49 horizontes das três séries exclui zero na banda de 90\%, e a faixa escura de 68\% exclui zero em 1 dos 147 pares, no painel do VIX em $h = 8$.}}
\label{fig:placebos}
\end{center}
\end{figure}

Um choque monetário brasileiro não deve mover ativos globais, e a banda de 90\% não o detecta movendo nenhum deles em nenhum horizonte. O que os placebos não descartam é um confundidor doméstico, porque uma surpresa fiscal brasileira deixaria igualmente parados o VIX, o MSCI de mercados emergentes e a incerteza de política econômica americana, e essa alternativa incide sobre a camada do instrumento.

\subsection{Invertibilidade do sistema}
\label{sec:invertibilidade}

A invertibilidade exige que os choques estruturais sejam recuperáveis por projeção linear sobre os valores correntes e passados das variáveis observadas \cite{stockwatson2018}. O SVAR-IV propaga o impacto identificado pelas inovações do VAR e, por isso, depende dessa condição. O LP-IV pode dispensá-la quando o instrumento satisfaz a condição lead-lag, mais restritiva. Na interpretação de variáveis omitidas dos mesmos autores, a invertibilidade implica que conhecer os choques verdadeiros não melhoraria as previsões do sistema. O DFM amplia o conjunto informacional e pode mitigar essa omissão, mas o aumento do número de variáveis não garante invertibilidade.

A não-causalidade de Granger do instrumento para as variáveis do VAR é uma implicação testável da invertibilidade, conforme a discussão da Seção 3 e a aplicação da Tabela 2 de \cite{stockwatson2018}. Em duas regressões sobre inovações fatoriais futuras, os valores-$p$ são $0{,}711$ com três liderações e $0{,}971$ com seis. O teste canônico inclui defasagens do instrumento nas cinco equações do VAR dos fatores e testa seus coeficientes conjuntamente, mantendo irrestritas as defasagens dos fatores. Com seis defasagens do instrumento, o menor valor-$p$ de \textit{bootstrap} entre as cinco equações é $0{,}324$, de modo que a não-causalidade de Granger não é rejeitada.

Essa não-rejeição verifica apenas uma condição necessária e não estabelece a invertibilidade, que permanece hipótese mantida do modelo. O teste também tem baixa potência, pois cada equação contém trinta regressores irrestritos e testa seis coeficientes em cento e quarenta e sete observações. O diagnóstico não fornece, portanto, uma prova de invertibilidade.

\subsection{O filtro de sinal seleciona risco soberano?}
\label{sec:confound}

O filtro de \cite{jarocinski2020} descarta os dias em que juros e ações sobem juntos e retém os de co-movimento negativo. Uma surpresa fiscal doméstica tem juros em alta, ações em queda e câmbio em depreciação, exatamente o padrão que o filtro retém como política monetária, de modo que a regra pode estar selecionando notícia de risco soberano em vez de choque de política. A objeção incide sobre o resultado central da Seção 4.2, porque se os dias retidos forem dias de notícia fiscal, câmbio e prêmio de risco se movem juntos por construção do instrumento e não por um canal de transmissão.

O teste roda na frequência diária, sobre a mesma janela de quarta a quinta-feira em que a surpresa é medida, com duas medidas de risco soberano. O índice $d$ identifica cada quinta-feira válida, de modo que $\Delta R_d$ e $x_d$ são variações medidas de quarta para quinta-feira, enquanto $t$ identifica os meses do instrumento agregado. A principal medida de risco é o CDS soberano de cinco anos, que cobre os 95 pares de dias das 95 reuniões e tem $0{,}5\%$ de variações exatamente nulas, contra 94 pares e $8{,}3\%$ do EMBI+. O EMBI+ é publicado com duas casas em pontos percentuais, e sua menor variação não nula exprimível é de 1 ponto-base contra $0{,}005$ do CDS, arredondamento que reduz a covariância medida e tornaria um resultado nulo descartável como erro de medida. A regra de leitura foi fixada sobre o EMBI+ antes que os números do CDS existissem, e o CDS é julgado pela mesma regra. A série diária do CDS correlaciona $0{,}9994$ com a série mensal do painel ao longo de 156 meses, ainda que venham de fornecedores distintos, o que torna o teste diário e a resposta mensal do CDS o mesmo objeto.

Contaminação exige que o dia retido carregue mais notícia de risco por unidade de surpresa do que uma quinta-feira comum. Formalizamos essa comparação por
\begin{equation}
\Delta R_d = \alpha + \beta x_d + \gamma D_d
+ \delta (x_d D_d) + u_d,
\label{eq:sovereign-contamination}
\end{equation}
em que $x_d$ é a surpresa e $D_d$ indica um dia retido pelo filtro. A contaminação conjecturada implica $\delta>0$, mas a estimativa de $\delta$ é negativa nas cinco medidas da janela do evento. Nas 503 quintas-feiras sem reunião do Copom a variação do CDS carrega a surpresa a $0{,}436$ (t = $4{,}35$, $R^2$ $0{,}218$), enquanto nos 62 dias retidos pelo filtro o coeficiente cai para $0{,}140$, cerca de um terço. No EMBI+ a razão é a mesma, $0{,}326$ contra $0{,}099$. A interação entre a surpresa e a indicadora dos dias retidos é $-0{,}191$ no CDS (valor-$p$ de $0{,}170$), $-0{,}182$ no EMBI+ ($0{,}092$) e $-0{,}036$ no câmbio ($0{,}068$).

O coeficiente nos 62 dias retidos não é zero, e essa é a primeira ressalva que o resultado carrega. Medido pelo CDS ele vale $0{,}140$ com valor-$p$ de \textit{bootstrap} de $0{,}002$, contra $0{,}099$ e $0{,}109$ do EMBI+, de modo que a afirmação sustentada é a de menos risco soberano que em um dia comum, e não a de ausência de risco. Os dias que o filtro rejeita carregam risco com sinal negativo e significativo ($-0{,}285$, valor-$p$ de $0{,}010$), enquanto o conjunto das 95 reuniões é nulo ($-0{,}004$), o que mostra que a regra separa dois regimes de sinal oposto em vez de descartar ruído.

Ortogonalizar a surpresa ao risco soberano contemporâneo é um limite inferior, porque política monetária pode mover o spread e a operação remove também essa resposta. A operação atinge apenas os valores da surpresa, enquanto a seleção dos 62 dias pela máscara de produção permanece fixa, de modo que ela examina o conteúdo do instrumento e não a regra que o constrói. A estatística $\xi_{mp}$ passa de $6{,}27$ na produção para $6{,}62$ quando a regressão auxiliar usa EMBI+ e câmbio, e para $7{,}58$ quando o CDS entra. Nessa última versão, o câmbio responde $0{,}155$ contra $0{,}158$ real por dólar na produção e o CDS abre 29,7 contra 32,5 pontos-base, com bandas de 90\% que excluem zero nas duas versões. A força permanece abaixo de 10 na amostra completa, de modo que o exercício preserva a direção das respostas sem eliminar a ressalva de instrumento fraco.

\subsection{O filtro de sinal seleciona notícias do FOMC?}
\label{sec:fomc}

A coincidência com decisões do FOMC é frequente entre os dias que o filtro de sinal retém. Ela ocorre em 24 dos 62 dias, equivalentes a $38{,}7\%$ do conjunto. Esses 24 dias concentram $35{,}5\%$ da soma das magnitudes absolutas das surpresas diárias retidas, $\sum_{d\in\mathcal R}|x_d|$, em que $x_d$ é a surpresa no dia $d$ e $\mathcal R$ contém os 62 dias selecionados pelo filtro. O grupo inclui 8 dos 20 dias de maior alavancagem. Em 2025, sete das oito reuniões do Copom coincidiram com uma decisão do FOMC. Essa exposição permite que uma notícia de política monetária americana contamine tanto a surpresa de juros quanto o retorno acionário usados pelo filtro.

O horário dos anúncios atenua apenas a primeira dessas duas ameaças. Nos 35 dias em que Copom e FOMC coincidem, a mediana do movimento absoluto da UST de dois anos é de $5{,}00$ pontos-base de terça para quarta e de $3{,}00$ pontos-base de quarta para quinta. Para o S\&P 500, porém, as medianas são de $0{,}54\%$ e $0{,}74\%$, respectivamente. O comunicado do FOMC ocorre antes do fechamento de quarta-feira, de modo que o argumento temporal cobre a perna de taxa, mas a reação acionária continua dentro da janela usada para construir o instrumento.

A regra de veredito foi fixada antes da estimação e classificaria como contaminação uma exposição maior à notícia americana nos dias selecionados. Os testes com interação usam
\begin{equation}
x_d = \alpha + b'U_d + \gamma D_d
+ \theta'(U_dD_d) + u_d,
\qquad
H_0\colon \theta=0,
\label{eq:fomc-contamination}
\end{equation}
em que $U_d=(\Delta UST2_d,r^{S\&P}_d)'$ reúne a variação da UST de dois anos e o retorno do S\&P 500. Conforme o teste, $D_d$ indica coincidência com uma decisão do FOMC ou seleção pelo filtro. Esse bloco contemporâneo não explica a surpresa nos 62 dias retidos, com valor-$p$ de \textit{bootstrap} de $0{,}458$. A interação do bloco com a coincidência FOMC tem valor-$p$ de $0{,}466$, enquanto a interação com a máscara do filtro tem $0{,}511$. Os três resultados levam ao veredito de que não se detecta confundimento pelo FOMC no bloco observado.

Os cortes da amostra apontam na mesma direção. Nos 35 dias de coincidência, o bloco americano explica $0{,}5\%$ da variância da surpresa, com $R^2=0{,}005$, enquanto nos 33 dias rejeitados pelo filtro explica $10{,}8\%$, com $R^2=0{,}108$. Se a coincidência injetasse notícia do Fed na surpresa retida, a relação deveria ser mais forte no primeiro conjunto. Os dados mostram o padrão oposto, pois a maior associação com a notícia americana está nos dias descartados.

Retirar o bloco americano apenas dos valores move $\xi_{mp}$ de $6{,}27$ para $6{,}31$. A seleção das reuniões pela máscara de produção permanece fixa, de modo que o exercício mede a sensibilidade dos valores da surpresa sem reclassificar os dias retidos. A bateria não detecta confundimento no bloco contemporâneo observado, mas não exclui spillovers que cheguem por variáveis fora desse bloco.

%---------------
\section{Concluding remarks}
%%---------------
Um choque contracionista de 50 pontos-base deprecia o real em $5,\!55\%$ no impacto, enquanto o EMBI+ abre 32,0 pontos-base e o CDS soberano de cinco anos abre 43,4 pontos-base. As três bandas de 90\% excluem zero, resultado incompatível com o canal convencional do diferencial de juros e compatível com transmissão via prêmio de risco. A estatística $\xi_{mp}=7{,}65$ na amostra completa, contudo, exige que as bandas convencionais sejam lidas com cautela.

A curva de juros acompanha o resultado no impacto, pois o repasse cresce com a maturidade até 117,0 pontos-base no vértice de cinco anos. A reversão posterior das taxas e do risco soberano coincide com o modo quase unitário do VAR dos fatores, de modo que não fornece confirmação independente do mecanismo. O resultado identificado com maior precisão é, portanto, o movimento conjunto de câmbio, risco soberano e curva no curto prazo.

%A curva de juros reforça essa leitura. O repasse do choque cresce com a maturidade até cinco anos, superando a amplificação que desenhos comparáveis encontram em economias desenvolvidas, e reverte no médio prazo em estreita sincronia com o câmbio e o prêmio de risco soberano, os cinco vértices, o EMBI+ e o CDS revertem de sinal em janelas próximas, entre o décimo e o vigésimo quarto mês. O câmbio, no entanto, não acompanha essa reversão: sua resposta perde significância estatística já a partir do quinto mês, sugerindo que o efeito sobre o risco soberano é mais persistente do que seu repasse ao câmbio.

Os demais resultados indicam queda da atividade no impacto e alta dos preços ao produtor no curto prazo, ambas com bandas de 90\% que excluem zero. O crédito agropecuário e o de transporte expandem inicialmente, enquanto a contração posterior ocorre apenas na banda de 68\% e depende da persistência do VAR dos fatores. No bloco acionário, somente o IFIX exclui zero a 90\%, em $h = 1$ e $h = 2$, de modo que o painel não sustenta uma conclusão geral para ações.

Em dados diários, \cite{goncalves2025} usam identificação via heterocedasticidade e encontram apreciação cambial e nenhuma resposta do CDS. A diferença pode vir da frequência dos dados ou da hipótese identificadora, mas as estimativas disponíveis não separam essas duas explicações. Essa limitação impede interpretar o contraste entre os estudos como teste direto de dominância fiscal.

Do ponto de vista de política econômica, nossos resultados indicam que, em economias emergentes sujeitas a fragilidades fiscais, a eficácia do canal cambial da política monetária não deve ser assumida automaticamente a partir do mecanismo convencional da paridade descoberta da taxa de juros. Nesses contextos, a credibilidade fiscal, e não apenas a magnitude do choque monetário, passa a condicionar a direção da resposta cambial. Mais amplamente, nossos resultados sugerem que a transmissão da política monetária deve ser compreendida como o resultado da interação entre política monetária e percepção de risco soberano, e não apenas como consequência do diferencial de juros. Essa perspectiva amplia a interpretação tradicional do canal cambial e, no caso brasileiro, oferece uma explicação alternativa para episódios em que aumentos da taxa básica de juros são acompanhados por depreciação cambial, um padrão que pode ser relevante para outras economias emergentes sujeitas a preocupações semelhantes de sustentabilidade fiscal.

%\section*{Funding}
%Jo\~ao F. Caldeira gratefully acknowledges support provided by CNPq under grant 309448/2022-0.

%%---------------
%% Bibliografia
%%---------------
%\section*{References}
\bibliographystyle{authordate3} % bibliography style
\bibliography{references} % References file

\section*{Appendix}\label{anexo-A}

\begin{ThreePartTable}
\begin{TableNotes}
\footnotesize
\item[*] Código de transformação com que a série entra na extração de fatores e é recuperada em unidades econômicas na função de impulso-resposta. 1 = nível. 4 = log-nível. 6 = retorno mensal, com a função de impulso-resposta em pontos percentuais do próprio retorno, sem acumulação. Os índices de mercado da B3 entram como retorno mensal. Os volumes nominais (base monetária, agregados, crédito, PIB e recolhimentos compulsórios) entram em logaritmo. As séries com sazonalidade detectada são ajustadas por X-13ARIMA-SEATS \cite{sax2018seasonal}. Todas as séries são padronizadas pelo desvio padrão da primeira diferença \cite{barigozzi2016non} antes da extração de fatores. As séries em nível não são diferenciadas, a não-estacionariedade é acomodada pela padronização.
\end{TableNotes}

\footnotesize
\setlength{\tabcolsep}{4pt}
\begin{longtable}{@{}p{5.2cm} l l c@{}}
\caption{Lista de Variáveis Utilizadas}\label{tab:lista_variaveis} \\
\toprule
\textbf{Variável} & \textbf{Fonte} & \textbf{Grupo} & \textbf{Código} \\
\midrule
\endfirsthead

\toprule
\textbf{Variável} & \textbf{Fonte} & \textbf{Grupo} & \textbf{Código} \\
\midrule
\endhead



\bottomrule
\insertTableNotes
\endlastfoot

Câmbio USD/BRL & Yahoo Finance & Câmbio & 1 \\
Câmbio EUR/BRL & Yahoo Finance & Câmbio & 1 \\
Câmbio CNY/BRL & Yahoo Finance & Câmbio & 1 \\
Câmbio ARS/BRL & Yahoo Finance & Câmbio & 1 \\
Câmbio INR/BRL & Yahoo Finance & Câmbio & 1 \\
Selic (meta acumulada) & BCB (4189) & Juros e curva & 1 \\
DI 3 meses & Svensson/B3 & Juros e curva & 1 \\
DI 6 meses & Svensson/B3 & Juros e curva & 1 \\
DI 1 ano & Svensson/B3 & Juros e curva & 1 \\
DI 2 anos & Svensson/B3 & Juros e curva & 1 \\
DI 5 anos & Svensson/B3 & Juros e curva & 1 \\
DI 10 anos & Svensson/B3 & Juros e curva & 1 \\
Base monetária & BCB (1788) & Agregados monetários & 4 \\
Meios de pagamento & BCB (27789) & Agregados monetários & 4 \\
Depósitos à vista & BCB (27790) & Agregados monetários & 4 \\
Depósitos de poupança & BCB (1835) & Agregados monetários & 4 \\
M1 & BCB (27791) & Agregados monetários & 4 \\
M2 & BCB (27810) & Agregados monetários & 4 \\
M3 & BCB (27813) & Agregados monetários & 4 \\
Crédito agropecuário & BCB (22027) & Crédito & 4 \\
Crédito indústria & BCB (22043) & Crédito & 4 \\
Crédito construção & BCB (22030) & Crédito & 4 \\
Crédito comércio & BCB (22036) & Crédito & 4 \\
Crédito transporte & BCB (22037) & Crédito & 4 \\
Crédito pessoa física & BCB (22050) & Crédito & 4 \\
Spread ICC pessoa jurídica & BCB (27444) & Crédito & 1 \\
Spread ICC pessoa física & BCB (27445) & Crédito & 1 \\
Saldo de crédito total & BCB (20542) & Crédito & 4 \\
Recolhimentos compulsórios & BCB (17633) & Crédito & 4 \\
Consumo de gasolina & BCB (1393) & Atividade econômica & 1 \\
Consumo de GLP & BCB (1394) & Atividade econômica & 1 \\
Consumo de óleo combustível & BCB (1395) & Atividade econômica & 1 \\
Consumo de óleo diesel & BCB (1396) & Atividade econômica & 1 \\
Consumo de demais derivados & BCB (1397) & Atividade econômica & 1 \\
Consumo de álcool & BCB (1401) & Atividade econômica & 1 \\
Utilização da capacidade instalada (CNI) & BCB (28554) & Atividade econômica & 1 \\
Vendas no varejo & BCB (1455) & Atividade econômica & 1 \\
Vendas de serviços & BCB (23982) & Atividade econômica & 1 \\
Produção de automóveis & BCB (7384) & Atividade econômica & 1 \\
Produção de veículos leves & BCB (7385) & Atividade econômica & 1 \\
Produção de caminhões & BCB (7386) & Atividade econômica & 1 \\
Produção de ônibus & BCB (7387) & Atividade econômica & 1 \\
Produção de bens de consumo & BCB (21865) & Atividade econômica & 1 \\
Produção de bens duráveis & BCB (21866) & Atividade econômica & 1 \\
Produção de bens não duráveis & BCB (21867) & Atividade econômica & 1 \\
Produção de bens de capital & BCB (21863) & Atividade econômica & 1 \\
Produção da indústria de transformação & BCB (21862) & Atividade econômica & 1 \\
Produção da indústria extrativa & BCB (21861) & Atividade econômica & 1 \\
Utilização da capacidade instalada (FGV) & BCB (24352) & Atividade econômica & 1 \\
Consumo de energia comercial & BCB (1402) & Atividade econômica & 1 \\
Consumo de energia residencial & BCB (1403) & Atividade econômica & 1 \\
Consumo de energia industrial & BCB (1404) & Atividade econômica & 1 \\
Consumo de energia outros & BCB (1405) & Atividade econômica & 1 \\
PIB acumulado no ano & BCB (4381) & Atividade econômica & 4 \\
IBC-Br & BCB (24364) & Atividade econômica & 1 \\
Índice de confiança do consumidor & BCB (4393) & Atividade econômica & 1 \\
Índice de confiança de serviços & BCB (17660) & Atividade econômica & 1 \\
Emprego formal Sudeste & BCB (13901) & Mercado de trabalho & 1 \\
Emprego formal Sul & BCB (13564) & Mercado de trabalho & 1 \\
Emprego formal Nordeste & BCB (13941) & Mercado de trabalho & 1 \\
Emprego formal Norte & BCB (28152) & Mercado de trabalho & 1 \\
Emprego formal Centro-Oeste & BCB (21991) & Mercado de trabalho & 1 \\
Salário mínimo & BCB (1619) & Mercado de trabalho & 1 \\
População na força de trabalho & BCB (24378) & Mercado de trabalho & 1 \\
População ocupada & BCB (28543) & Mercado de trabalho & 1 \\
Taxa de desemprego & BCB (24369) & Mercado de trabalho & 1 \\
Horas trabalhadas na indústria & BCB (28556) & Mercado de trabalho & 1 \\
Tempo de procura: até 1 mês & IBGE/SIDRA (1616) & Mercado de trabalho & 1 \\
Tempo de procura: 1 mês a 1 ano & IBGE/SIDRA (1616) & Mercado de trabalho & 1 \\
Tempo de procura: 1 a 2 anos & IBGE/SIDRA (1616) & Mercado de trabalho & 1 \\
Tempo de procura: 2 anos ou mais & IBGE/SIDRA (1616) & Mercado de trabalho & 1 \\
IPCA & BCB (433) & Inflação & 1 \\
IPCA índice de difusão & BCB (21379) & Inflação & 1 \\
IPCA núcleo EX0 & BCB (29677) & Inflação & 1 \\
IPCA núcleo EX1 & BCB (29678) & Inflação & 1 \\
IPCA núcleo médias aparadas & BCB (16122) & Inflação & 1 \\
IGP-M & BCB (189) & Inflação & 1 \\
IPC & BCB (191) & Inflação & 1 \\
IPC núcleo & BCB (4467) & Inflação & 1 \\
INCC & BCB (192) & Inflação & 1 \\
INPC & BCB (188) & Inflação & 1 \\
IPP & IPEA (IPP12\_IPPCG12) & Inflação & 1 \\
Índice de commodities agropecuária & BCB (27575) & Commodities & 1 \\
Índice de commodities metal & BCB (27576) & Commodities & 1 \\
Índice de commodities energia & BCB (27577) & Commodities & 1 \\
Ibovespa (retorno mensal) & B3 & Índices de mercado & 6 \\
IDIV (retorno mensal) & B3 & Índices de mercado & 6 \\
IFIX (retorno mensal) & B3 & Índices de mercado & 6 \\
IFNC (retorno mensal) & B3 & Índices de mercado & 6 \\
IMAT (retorno mensal) & B3 & Índices de mercado & 6 \\
IMOB (retorno mensal) & B3 & Índices de mercado & 6 \\
SMLL (retorno mensal) & B3 & Índices de mercado & 6 \\
EMBI+ Brasil & JP Morgan & Risco e incerteza & 1 \\
CDS soberano 5 anos & Investing.com & Risco e incerteza & 1 \\
MSCI mercados emergentes & Investing.com & Risco e incerteza & 1 \\
VIX & CBOE & Risco e incerteza & 1 \\
EPU Brasil & policyuncertainty.com & Risco e incerteza & 1 \\
EPU Canadá & policyuncertainty.com & Risco e incerteza & 1 \\
EPU Chile & policyuncertainty.com & Risco e incerteza & 1 \\
EPU China & policyuncertainty.com & Risco e incerteza & 1 \\
EPU Alemanha & policyuncertainty.com & Risco e incerteza & 1 \\
EPU Índia & policyuncertainty.com & Risco e incerteza & 1 \\
EPU Rússia & policyuncertainty.com & Risco e incerteza & 1 \\
EPU Estados Unidos & policyuncertainty.com & Risco e incerteza & 1 \\
Dívida bruta do governo geral (\% PIB) & BCB (13762) & Fiscal & 1 \\
Dívida líquida do setor público (\% PIB) & BCB (4513) & Fiscal & 1 \\
Resultado primário do setor público & BCB (4649) & Fiscal & 1 \\
Focus: IPCA suavizado 12 meses & BCB/Focus & Expectativas & 1 \\
Focus: Selic no fim do ano seguinte & BCB/Focus & Expectativas & 1 \\
Focus: PIB no ano seguinte & BCB/Focus & Expectativas & 1 \\
Focus: câmbio no fim do ano seguinte & BCB/Focus & Expectativas & 1 \\

\end{longtable}
\end{ThreePartTable}





\end{document}
